%%%%%%
%
% $Autor: Wings $
% $Datum: 2019-03-05 08:03:15Z $
% $Pfad: Elektromagnet $
% $Version: 4250 $
% !TeX spellcheck = en_GB/de_DE
% !TeX encoding = utf8
% !TeX root = filename 
% !TeX TXS-program:bibliography = txs:///biber
%
%%%%%%


\chapter{Elektromagnet}


\section{Specification}

Kurzbeschreibung:



DC 4.5V 0.11A 10g Fähigkeit 3mm Hub Ziehen Typ Magnet Elektromagnet

\medskip

\begin{tabular}{ll}
	Marke             & Uxcell \\
	Herstellernummer  & nicht zutreffend \\
	upc               & 702105986189 \\	
	Produktname       & DC Magnet Elektromagnet \\
	Nennspannung      & DC 4.5V \\
	Typ	              & Ziehen \\
	Nennstrom	      & 0.11A \\
	Leistung          & 0.5W \\
	Fähigkeit and Hub & 10g/3mm \\
	Angetriebene Rate & 50\% \\
	Körpergröße	      & 20 $\times$ 12 $\times$ 11mm / 0,78'' $\times$ 0,47'' $\times$ 0,43' (L $\cdot$  B $\cdot$ D) \\
	Stößel Stab Größe & 4 $\times$ 8mm / 0.16'' $\times$ 0.3'' (D $\cdot$ L) \\
	Kabellänge        & 10cm/3.9'' \\
	Material          & Metall, Elektronische Teile \\
	Hauptfarbe        & siehe Abbildung \\
	Gewicht           & 15g \\
	Lieferumfang      & 1 $\times$ DC Solenoid Elektromagnet \\
\end{tabular}


	
\section{Beschreibung}
	
DC Solemoid Elektromagneten hauptsächlich in Verkaufsautomaten, Transportausrüstung, Büroausstattung Haushaltsgerät, mechanisch usw. verwendet.


Wenn der Elektromagnet erregt wird, erfolgt das Ziehen des Kolbens.
	

\begin{center}	
	
	\includegraphics[width=0.8\textwidth]{Elektromagnet/Uxcell702105986189}
	
	\captionof{figure}{Elektromagnet}
	
\end{center}
	
\medskip

Bezug über Ebay:	\URL{https://www.ebay.de/itm/134941988659}
	
	
\section{Wasserpumpe}
	
	
	\URL{https://www.ebay.de/itm/296892083454}
	
	\URL{https://www.ebay.de/itm/353234654746}
	
	
\section{Unterschiede der Arduino Nano BLE Sense}
	
	
	
	Arduino Nano 33 BLE Sense Lite hat  keine HTS221-Temperatur- und Feuchtigkeitssensoren, sondern nur den LPS22HB-Drucksensor, der allerdings auch die Temperatur messen kann
	
	Arduino Nano 33 BLE Sense Rev2 hat  einen HTS221-Temperatur- und Feuchtigkeitssensor und den LPS22HB-Drucksensor, der allerdings auch die Temperatur messen kann.
	
	
	Unterschiede:
	
	Inertiale Messeinheit (IMU):
	Rev 1: Verfügt über eine einzelne 9-Achsen-IMU.
	Rev 2: Hat eine Kombination aus zwei IMUs: eine 6-Achsen-IMU (BMI270) und eine 3-Achsen-IMU (BMM150), was die Möglichkeiten zur Bewegungserkennung erweitert1.
	
	Design und Zugänglichkeit:
	Rev 2: Integriert neue Pads und Testpunkte für USB, SWDIO und SWCLK, was den Zugang zu diesen wichtigen Punkten auf der Platine erleichtert1.
	Rev 1: Hat diese zusätzlichen Pads und Testpunkte nicht.
	Stromversorgung:
	
	Rev 2: Verwendet den MP2322 als Stromversorgungskomponente, was die Leistung verbessert1.
	
	Rev 1: Nutzt eine andere Stromversorgungskomponente.
	Der Arduino Nano 33 BLE Sense Rev 1 verwendet den MPS MP2144 als Stromversorgungskomponente
	
	


